\documentclass[10pt,a4paper]{amsart}
\usepackage{amsmath,amssymb,amsthm,graphicx}
\usepackage{verbatim}
\newenvironment{reference}{\comment}{\endcomment}
%\newenvironment{reference}{}{}
\newenvironment{slogan}{\comment}{\endcomment}
\newenvironment{history}{\comment}{\endcomment}
% For commutative diagrams we use Xy-pic
\usepackage[all]{xy}

% We use 2cell for 2-commutative diagrams.
\xyoption{2cell}
% Theorem environments.
%
\theoremstyle{plain}
\newtheorem{theorem}[subsection]{Theorem}
\newtheorem{proposition}[subsection]{Proposition}
\newtheorem{lemma}[subsection]{Lemma}

\theoremstyle{definition}
\newtheorem{definition}[subsection]{Definition}
\newtheorem{example}[subsection]{Example}
\newtheorem{exercise}[subsection]{Exercise}
\newtheorem{situation}[subsection]{Situation}

\theoremstyle{remark}
\newtheorem{remark}[subsection]{Remark}
\newtheorem{remarks}[subsection]{Remarks}

\numberwithin{equation}{subsection}

% Macros
%
\def\lim{\mathop{\mathrm{lim}}\nolimits}
\def\colim{\mathop{\mathrm{colim}}\nolimits}
\def\Spec{\mathop{\mathrm{Spec}}}
\def\Hom{\mathop{\mathrm{Hom}}\nolimits}
\def\Ext{\mathop{\mathrm{Ext}}\nolimits}
\def\SheafHom{\mathop{\mathcal{H}\!\mathit{om}}\nolimits}
\def\SheafExt{\mathop{\mathcal{E}\!\mathit{xt}}\nolimits}
\def\Sch{\mathit{Sch}}
\def\Mor{\mathop{\mathrm{Mor}}\nolimits}
\def\Ob{\mathop{\mathrm{Ob}}\nolimits}
\def\Sh{\mathop{\mathit{Sh}}\nolimits}
\def\NL{\mathop{N\!L}\nolimits}
\def\CH{\mathop{\mathrm{CH}}\nolimits}
\def\proetale{{pro\text{-}\acute{e}tale}}
\def\etale{{\acute{e}tale}}
\def\QCoh{\mathit{QCoh}}
\def\Ker{\mathop{\mathrm{Ker}}}
\def\Im{\mathop{\mathrm{Im}}}
\def\Coker{\mathop{\mathrm{Coker}}}
\def\Coim{\mathop{\mathrm{Coim}}}

% Boxtimes
%
\DeclareMathSymbol{\boxtimes}{\mathbin}{AMSa}{"02}

%
% Macros for moduli stacks/spaces
%
\def\QCohstack{\mathcal{QC}\!\mathit{oh}}
\def\Cohstack{\mathcal{C}\!\mathit{oh}}
\def\Spacesstack{\mathcal{S}\!\mathit{paces}}
\def\Quotfunctor{\mathrm{Quot}}
\def\Hilbfunctor{\mathrm{Hilb}}
\def\Curvesstack{\mathcal{C}\!\mathit{urves}}
\def\Polarizedstack{\mathcal{P}\!\mathit{olarized}}
\def\Complexesstack{\mathcal{C}\!\mathit{omplexes}}
% \Pic is the operator that assigns to X its picard group, usage \Pic(X)
% \Picardstack_{X/B} denotes the Picard stack of X over B
% \Picardfunctor_{X/B} denotes the Picard functor of X over B
\def\Pic{\mathop{\mathrm{Pic}}\nolimits}
\def\Picardstack{\mathcal{P}\!\mathit{ic}}
\def\Picardfunctor{\mathrm{Pic}}
\def\Deformationcategory{\mathcal{D}\!\mathit{ef}}
\usepackage{cleveref}
\crefname{originallemma}{original lemma}{original lemmas}
\Crefname{originallemma}{Original lemma}{Original lemmas}
\makeatletter
\newlabel{coherent-lemma-quasi-coherent-affine-cohomology-zero}{{01XB}{source}{}{sourcepointer-1}{}}
\newlabel{coherent-lemma-quasi-coherent-affine-cohomology-zero@cref}{{[originallemma][0][]01XB}{[1][0][]source}{}{}{}}
\makeatother
\crefname{originallemma}{original lemma}{original lemmas}
\Crefname{originallemma}{Original lemma}{Original lemmas}
\makeatletter
\newlabel{lemma-higher-direct-image}{{0720}{source}{}{sourcepointer-2}{}}
\newlabel{lemma-higher-direct-image@cref}{{[originallemma][0][]0720}{[1][0][]source}{}{}{}}
\makeatother
\crefname{originallemma}{original lemma}{original lemmas}
\Crefname{originallemma}{Original lemma}{Original lemmas}
\makeatletter
\newlabel{sites-cohomology-lemma-Leray}{{0732}{source}{}{sourcepointer-3}{}}
\newlabel{sites-cohomology-lemma-Leray@cref}{{[originallemma][0][]0732}{[1][0][]source}{}{}{}}
\makeatother
\crefname{originallemma}{original lemma}{original lemmas}
\Crefname{originallemma}{Original lemma}{Original lemmas}
\makeatletter
\newlabel{spaces-properties-lemma-pushforward-etale-base-change-modules}{{03LX}{source}{}{sourcepointer-4}{}}
\newlabel{spaces-properties-lemma-pushforward-etale-base-change-modules@cref}{{[originallemma][0][]03LX}{[1][0][]source}{}{}{}}
\makeatother
\title{Cohomology over an affine etale neighbourhood is the global sections of the higher direct image}
\author{The Stacks Project Authors}
\date{stacks-project:spaces-cohomology.tex:local-snapshot-20261012; GNU FDL 1.2 or later}
\begin{document}
\maketitle
\noindent\textit{Editorial note.} All delivered mathematical text is the author's own, copied from the pinned Stacks Project chapter file spaces-cohomology.tex; no statement or proof step was rewritten, reordered or removed.\par\medskip
\noindent\textit{Editorial note.} Preamble provenance: the pinned chapter file opens with \textbackslash{}input\{preamble\} and carries no preamble text of its own, so the project's single Stacks-wide preamble staged at sources/stacks-spaces-cohomology/preamble.tex is excerpted verbatim as three spans, lines 8-12 (verbatim package and the project's own reference/slogan/history display environments), lines 14-18 (Xy-pic setup) and lines 154-217 (theorem declarations and mathematical macros). The wrapper class and the amsmath/amssymb/amsthm/graphicx base packages are supplied by the assembler. No macro was invented or redefined.\par\medskip
\noindent\textit{Editorial note.} Same-chapter pointers are located at their \textbackslash{}label lines inside the pinned chapter file; cross-chapter pointers are located in the project's own tag list staged in the snapshot (sources/stacks-chow/tags.txt), at the verbatim lines carrying their canonical tags.\par\medskip
\noindent\textit{Allowed standard prerequisites (editorial).}\begin{itemize}
\item The higher direct images commute with etale localization. Properties of Spaces, label spaces-properties-lemma-pushforward-etale-base-change-modules (invoked at proof line 468). It reduces the statement for an arbitrary \$Y\$ to the case where \$Y = V\$ is affine.
\item The Leray spectral sequence of a morphism of algebraic spaces. Cohomology on Sites, label sites-cohomology-lemma-Leray (invoked at proof line 474). It gives the spectral sequence \$E\_2\textasciicircum{}\{p,q\} = H\textasciicircum{}p(Y, R\textasciicircum{}qf\_*\textbackslash{}mathcal\{F\})\$ converging to \$H\textasciicircum{}\{p+q\}(X, \textbackslash{}mathcal\{F\})\$.
\item Quasi-coherence of the higher direct images along a quasi-compact quasi-separated morphism. Same chapter spaces-cohomology.tex, label lemma-higher-direct-image (invoked at proof line 475). It puts \$R\textasciicircum{}qf\_*\textbackslash{}mathcal\{F\}\$ inside the class of quasi-coherent modules on the affine base, whose higher cohomology vanishes.
\item Vanishing of the higher cohomology of quasi-coherent modules on an affine scheme. Cohomology of Schemes, label coherent-lemma-quasi-coherent-affine-cohomology-zero (invoked at proof line 478). It kills the terms \$E\_2\textasciicircum{}\{p,q\}\$ with \$p > 0\$, so the spectral sequence degenerates at \$E\_2\$ and the claim follows.
\end{itemize}\medskip
\ifdefined\hypertarget\hypertarget{sourcepointer-1}{}\fi
\noindent\textit{Source reference.} Original lemma 01XB: sources/stacks-chow/tags.txt, source lines 2512--2512. Navigation pointer for \textbackslash{}ref\{coherent-lemma-quasi-coherent-affine-cohomology-zero\} used inside the counted span (line 478); the argument that consumes it is retained in the span.\par\medskip
\ifdefined\hypertarget\hypertarget{sourcepointer-2}{}\fi
\noindent\textit{Source reference.} Original lemma 0720: sources/stacks-spaces-cohomology/pinned-original-source.tex, source lines 339--339. Navigation pointer for \textbackslash{}ref\{lemma-higher-direct-image\} used inside the counted span (line 475); the argument that consumes it is retained in the span.\par\medskip
\ifdefined\hypertarget\hypertarget{sourcepointer-3}{}\fi
\noindent\textit{Source reference.} Original lemma 0732: sources/stacks-chow/tags.txt, source lines 9046--9046. Navigation pointer for \textbackslash{}ref\{sites-cohomology-lemma-Leray\} used inside the counted span (line 474); the argument that consumes it is retained in the span.\par\medskip
\ifdefined\hypertarget\hypertarget{sourcepointer-4}{}\fi
\noindent\textit{Source reference.} Original lemma 03LX: sources/stacks-chow/tags.txt, source lines 4729--4729. Navigation pointer for \textbackslash{}ref\{spaces-properties-lemma-pushforward-etale-base-change-modules\} used inside the counted span (line 469); the argument that consumes it is retained in the span.\par\medskip
\begin{lemma}
\label{lemma-quasi-coherence-higher-direct-images-application}
Let $S$ be a scheme. Let $f : X \to Y$ be a quasi-separated and quasi-compact
morphism of algebraic spaces over $S$. For any quasi-coherent
$\mathcal{O}_X$-module $\mathcal{F}$ and any affine object $V$ of
$Y_\etale$ we have
$$
H^q(V \times_Y X, \mathcal{F}) = H^0(V, R^qf_*\mathcal{F})
$$
for all $q \in \mathbf{Z}$.
\end{lemma}

\begin{proof}
Since formation of $Rf_*$ commutes with \'etale localization
(Properties of Spaces, Lemma
\ref{spaces-properties-lemma-pushforward-etale-base-change-modules})
we may replace $Y$ by $V$ and assume $Y = V$ is affine.
Consider the Leray spectral sequence
$E_2^{p, q} = H^p(Y, R^qf_*\mathcal{F})$
converging to $H^{p + q}(X, \mathcal{F})$, see
Cohomology on Sites, Lemma \ref{sites-cohomology-lemma-Leray}.
By Lemma \ref{lemma-higher-direct-image}
we see that the sheaves $R^qf_*\mathcal{F}$ are quasi-coherent. By
Cohomology of Schemes, Lemma
\ref{coherent-lemma-quasi-coherent-affine-cohomology-zero}
we see that $E_2^{p, q} = 0$ when $p > 0$.
Hence the spectral sequence degenerates at $E_2$ and we win.
\end{proof}

\end{document}
